% Options for packages loaded elsewhere
\PassOptionsToPackage{unicode}{hyperref}
\PassOptionsToPackage{hyphens}{url}
\documentclass[
  english,
  11pt,
]{article}
\usepackage{xcolor}
\usepackage[margin=1in]{geometry}
\usepackage{amsmath,amssymb}
\setcounter{secnumdepth}{5}
\usepackage{iftex}
\ifPDFTeX
  \usepackage[T1]{fontenc}
  \usepackage[utf8]{inputenc}
  \usepackage{textcomp} % provide euro and other symbols
\else % if luatex or xetex
  \usepackage{unicode-math} % this also loads fontspec
  \defaultfontfeatures{Scale=MatchLowercase}
  \defaultfontfeatures[\rmfamily]{Ligatures=TeX,Scale=1}
\fi
\usepackage{lmodern}
\ifPDFTeX\else
  % xetex/luatex font selection
\fi
% Use upquote if available, for straight quotes in verbatim environments
\IfFileExists{upquote.sty}{\usepackage{upquote}}{}
\IfFileExists{microtype.sty}{% use microtype if available
  \usepackage[]{microtype}
  \UseMicrotypeSet[protrusion]{basicmath} % disable protrusion for tt fonts
}{}
\makeatletter
\@ifundefined{KOMAClassName}{% if non-KOMA class
  \IfFileExists{parskip.sty}{%
    \usepackage{parskip}
  }{% else
    \setlength{\parindent}{0pt}
    \setlength{\parskip}{6pt plus 2pt minus 1pt}}
}{% if KOMA class
  \KOMAoptions{parskip=half}}
\makeatother
\usepackage{color}
\usepackage{fancyvrb}
\newcommand{\VerbBar}{|}
\newcommand{\VERB}{\Verb[commandchars=\\\{\}]}
\DefineVerbatimEnvironment{Highlighting}{Verbatim}{commandchars=\\\{\}}
% Add ',fontsize=\small' for more characters per line
\newenvironment{Shaded}{}{}
\newcommand{\AlertTok}[1]{\textcolor[rgb]{1.00,0.00,0.00}{\textbf{#1}}}
\newcommand{\AnnotationTok}[1]{\textcolor[rgb]{0.38,0.63,0.69}{\textbf{\textit{#1}}}}
\newcommand{\AttributeTok}[1]{\textcolor[rgb]{0.49,0.56,0.16}{#1}}
\newcommand{\BaseNTok}[1]{\textcolor[rgb]{0.25,0.63,0.44}{#1}}
\newcommand{\BuiltInTok}[1]{\textcolor[rgb]{0.00,0.50,0.00}{#1}}
\newcommand{\CharTok}[1]{\textcolor[rgb]{0.25,0.44,0.63}{#1}}
\newcommand{\CommentTok}[1]{\textcolor[rgb]{0.38,0.63,0.69}{\textit{#1}}}
\newcommand{\CommentVarTok}[1]{\textcolor[rgb]{0.38,0.63,0.69}{\textbf{\textit{#1}}}}
\newcommand{\ConstantTok}[1]{\textcolor[rgb]{0.53,0.00,0.00}{#1}}
\newcommand{\ControlFlowTok}[1]{\textcolor[rgb]{0.00,0.44,0.13}{\textbf{#1}}}
\newcommand{\DataTypeTok}[1]{\textcolor[rgb]{0.56,0.13,0.00}{#1}}
\newcommand{\DecValTok}[1]{\textcolor[rgb]{0.25,0.63,0.44}{#1}}
\newcommand{\DocumentationTok}[1]{\textcolor[rgb]{0.73,0.13,0.13}{\textit{#1}}}
\newcommand{\ErrorTok}[1]{\textcolor[rgb]{1.00,0.00,0.00}{\textbf{#1}}}
\newcommand{\ExtensionTok}[1]{#1}
\newcommand{\FloatTok}[1]{\textcolor[rgb]{0.25,0.63,0.44}{#1}}
\newcommand{\FunctionTok}[1]{\textcolor[rgb]{0.02,0.16,0.49}{#1}}
\newcommand{\ImportTok}[1]{\textcolor[rgb]{0.00,0.50,0.00}{\textbf{#1}}}
\newcommand{\InformationTok}[1]{\textcolor[rgb]{0.38,0.63,0.69}{\textbf{\textit{#1}}}}
\newcommand{\KeywordTok}[1]{\textcolor[rgb]{0.00,0.44,0.13}{\textbf{#1}}}
\newcommand{\NormalTok}[1]{#1}
\newcommand{\OperatorTok}[1]{\textcolor[rgb]{0.40,0.40,0.40}{#1}}
\newcommand{\OtherTok}[1]{\textcolor[rgb]{0.00,0.44,0.13}{#1}}
\newcommand{\PreprocessorTok}[1]{\textcolor[rgb]{0.74,0.48,0.00}{#1}}
\newcommand{\RegionMarkerTok}[1]{#1}
\newcommand{\SpecialCharTok}[1]{\textcolor[rgb]{0.25,0.44,0.63}{#1}}
\newcommand{\SpecialStringTok}[1]{\textcolor[rgb]{0.73,0.40,0.53}{#1}}
\newcommand{\StringTok}[1]{\textcolor[rgb]{0.25,0.44,0.63}{#1}}
\newcommand{\VariableTok}[1]{\textcolor[rgb]{0.10,0.09,0.49}{#1}}
\newcommand{\VerbatimStringTok}[1]{\textcolor[rgb]{0.25,0.44,0.63}{#1}}
\newcommand{\WarningTok}[1]{\textcolor[rgb]{0.38,0.63,0.69}{\textbf{\textit{#1}}}}
\usepackage{longtable,booktabs,array}
\newcounter{none} % for unnumbered tables
\usepackage{calc} % for calculating minipage widths
% Correct order of tables after \paragraph or \subparagraph
\usepackage{etoolbox}
\makeatletter
\patchcmd\longtable{\par}{\if@noskipsec\mbox{}\fi\par}{}{}
\makeatother
% Allow footnotes in longtable head/foot
\IfFileExists{footnotehyper.sty}{\usepackage{footnotehyper}}{\usepackage{footnote}}
\makesavenoteenv{longtable}
% definitions for citeproc citations
\NewDocumentCommand\citeproctext{}{}
\NewDocumentCommand\citeproc{mm}{%
  \begingroup\def\citeproctext{#2}\cite{#1}\endgroup}
\makeatletter
 % allow citations to break across lines
 \let\@cite@ofmt\@firstofone
 % avoid brackets around text for \cite:
 \def\@biblabel#1{}
 \def\@cite#1#2{{#1\if@tempswa , #2\fi}}
\makeatother
\newlength{\cslhangindent}
\setlength{\cslhangindent}{1.5em}
\newlength{\csllabelwidth}
\setlength{\csllabelwidth}{3em}
\newenvironment{CSLReferences}[2] % #1 hanging-indent, #2 entry-spacing
 {\begin{list}{}{%
  \setlength{\itemindent}{0pt}
  \setlength{\leftmargin}{0pt}
  \setlength{\parsep}{0pt}
  % turn on hanging indent if param 1 is 1
  \ifodd #1
   \setlength{\leftmargin}{\cslhangindent}
   \setlength{\itemindent}{-1\cslhangindent}
  \fi
  % set entry spacing
  \setlength{\itemsep}{#2\baselineskip}}}
 {\end{list}}
\usepackage{calc}
\newcommand{\CSLBlock}[1]{\hfill\break\parbox[t]{\linewidth}{\strut\ignorespaces#1\strut}}
\newcommand{\CSLLeftMargin}[1]{\parbox[t]{\csllabelwidth}{\strut#1\strut}}
\newcommand{\CSLRightInline}[1]{\parbox[t]{\linewidth - \csllabelwidth}{\strut#1\strut}}
\newcommand{\CSLIndent}[1]{\hspace{\cslhangindent}#1}
\ifLuaTeX
\usepackage[bidi=basic,shorthands=off]{babel}
\else
\usepackage[bidi=default,shorthands=off]{babel}
\fi
\ifLuaTeX
  \usepackage{selnolig} % disable illegal ligatures
\fi
\setlength{\emergencystretch}{3em} % prevent overfull lines
\providecommand{\tightlist}{%
  \setlength{\itemsep}{0pt}\setlength{\parskip}{0pt}}
\usepackage{bookmark}
\IfFileExists{xurl.sty}{\usepackage{xurl}}{} % add URL line breaks if available
\urlstyle{same}
\hypersetup{
  pdftitle={Astrocyte: Experience Reuse and Collective Adaptation in Decentralized Agent Swarms},
  pdfauthor={Authors: TO BE PROVIDED},
  pdflang={en},
  hidelinks,
  pdfcreator={LaTeX via pandoc}}

\title{Astrocyte: Experience Reuse and Collective Adaptation in
Decentralized Agent Swarms}
\author{Authors: TO BE PROVIDED}
\date{19 October 2026 --- method and protocol draft 0.4}

\begin{document}
\maketitle
\begin{abstract}
Reusable experience can influence how agents solve tasks and choose
collaborators, while collaboration generates evidence that can revise
both experience and future relationships. We propose Astrocyte as a
framework for studying this coupled adaptation with fixed
foundation-model parameters and constrained execution budgets. The
proposed method combines versioned task and collaboration experience,
condition-specific directed relationship scores, evidence-gated updates,
and a shared decision rule that can be executed centrally or by
individual members. A three-factor, eight-configuration protocol
separates experience updates, relationship updates, and decision
location. Its primary estimand is the interaction between the two
updates under local decisions; secondary estimands measure how this
interaction and the total update gain vary with decision location.
Complete task sequences form paired replication units, and execution
feedback is separated from sealed evaluation. We specify cost
accounting, member perturbations, and a prospective system-iteration
case. This manuscript presents a method and an unregistered experimental
protocol. Current implementation evidence covers selected product
components; the research policy, complete experimental controls, and
formal measurements remain incomplete. No formal experimental results
are reported, and no positive adaptation or decentralization effect is
assumed.
\end{abstract}

\textbf{Affiliations: TO BE PROVIDED. ORCID identifiers: TO BE
PROVIDED.}

\textbf{Status.} This is a complete first manuscript for review, based
on the project positioning document v3 and experimental preregistration
draft v0.1, both dated 19 October 2026. The protocol has not been
formally registered. Model versions, tasks, budgets, sample sizes, and
several analysis details require confirmation before registration.
Nothing in this draft constitutes an empirical performance claim for the
proposed method.

\section{Introduction}\label{introduction}

An agent can preserve a useful procedure after completing a task.
Reusing that procedure may change which task it attempts, which
collaborator it requests, or which check it performs. Those choices
change the evidence available for subsequent learning. A collaboration
network can similarly favor partners whose recent contributions have
worked under relevant conditions. These processes form a coupled
feedback system: experience affects collaboration, and collaboration
affects the production and revision of experience.

The scientific question is whether updating both processes yields a
practically useful interaction over continuous tasks. A system might
benefit from each update independently and still show little additional
value from their combination. It might also exhibit a negative
interaction when an erroneous experience repeatedly directs work toward
an unsuitable partner. Measuring only the full system against a static
baseline cannot distinguish these possibilities.

We therefore propose a factorial study of experience updates,
relationship updates, and the location of collaboration decisions.
Foundation-model parameters remain fixed within the main study. The
external experience state and condition-specific relationship records
may change according to the assigned configuration. The word
\emph{decentralized} refers to members submitting their own task and
partner decisions. Authorization, task claims, durable records, and
budget enforcement may use shared services. Each Mission retains one
authoritative execution backend.

The proposed research focuses on connecting exact experience versions to
checkable collaboration evidence, defining condition-specific
relationship updates, and separating update interaction from
deployment-dependent decision-location effects. This manuscript provides
a method specification, a proposed analysis procedure, and an
implementation audit. Their substantive originality still requires
assessment against closely related methods, and the incomplete research
implementation requires controlled validation. Factorial contrast
algebra and general statistical procedures are established tools; they
are not presented as independent theoretical breakthroughs. Benefits
remain unmeasured.

The main study is distinct from the proposed system-iteration case. In
the main study, the code and decision rule remain fixed across
conditions. In the case, selected code or configuration changes are
themselves prospective research artifacts. Mixing these settings would
confound adaptation of external state with changes to the underlying
agent system.

\section{Related work and the scope of the research
question}\label{related-work-and-the-scope-of-the-research-question}

Strategy Genes studies compact, structured experience representations
and selective incorporation of failure information with fixed model
parameters (Wang et al. 2026, secs. 3--4 and Appendix A). This motivates
controlling experience representation and its cost. Astrocyte's proposed
factorial question concerns the joint effect of updating that experience
and persistent collaboration relationships.

ShareMem separates transferable experience from user-bound preferences
and studies two-stage consolidation and scoped retrieval (Hu et al.
2026, sec. 3 and Appendices A--B). It motivates recording coverage,
provenance, and construction costs. Its results do not determine how
continuous relationship updates interact with experience updates in the
software-task setting proposed here.

EDV separates heterogeneous execution, third-party distillation, and
executor consensus verification (Zhu et al. 2026, sec. 4 and Appendix
G). Consensus remains a model judgment with possible shared bias.
Astrocyte therefore proposes separate records for model agreement and
independently checkable execution evidence.

MAAGL already uses historical trajectories for confidence estimation,
collaborator selection, and experience learning in graph tasks (Qu et
al. 2026, secs. IV-C--IV-D). Historical experience guiding partner
selection is consequently an existing mechanism. The proposed
distinction here is a controlled interaction study with
condition-specific directed relationship updates over sequential
software tasks. This is a research target, with novelty and
effectiveness still requiring assessment. The cited author list follows
the paper's HTML and PDF title pages; the abstract-page indexing
discrepancy remains documented in the source audit.

DecentMem maintains agent-private exploitation and exploration memory
pools, reuses collaboration trajectories, and updates online routing
from stage-wise LLM judgments (Hao et al. 2026, secs. 3--4).
Decentralized memory combined with collaboration experience is therefore
also established prior work. Our proposed decision-location comparison
gives central and local decision makers the same designated information
view; it does not independently manipulate private versus shared memory
architecture. The theoretical claims in DecentMem have their own
search-space and reward assumptions, which cannot be transferred to the
changing authorized candidate sets and discounted evidence proposed
here. Neither a generic memory combination nor the word decentralization
establishes this draft's novelty.

Market-based post-training jointly trains action and bidding policies
and evaluates population changes (Xu et al. 2026, sec. 2 and 3.4).
Astrocyte proposes external-state adaptation with fixed model parameters
and separately manipulates decision location. Member perturbations are a
relevant evaluation dimension, with their effects to be measured under
the present protocol.

Darwin Gödel Machine maintains an archive of code-modified agents and
evaluates candidate changes (Zhang et al. 2025, secs. 3--4 and Appendix
F). It motivates the prospective iteration case and the need to protect
evaluation from candidate-controlled instrumentation. EvoMap's official
Swarm documentation describes shared memory and history-sensitive
collaboration (EvoMap, n.d., ``Shared Memory'' and ``Intelligent
Orchestration''). That documentation provides design context; it
supplies no causal estimate for Astrocyte.

The reviewed sources establish substantial prior work on experience
reuse and adaptive coordination. This draft makes no claim to an
exhaustive literature review or priority for combining memory with
collaborator selection. The intended evidence would concern the
specified update interaction and its limits.

\section{Setting and state}\label{setting-and-state}

Let a task sequence contain \(T\) scheduled tasks in a fixed order. A
member has an identity, a known capability version, a current
availability state, and an authorization scope. A decision \(d\) is
associated with scheduled position \(t\) and an authoritative
within-position order. Its recorded input contains only information
available before submission. Foundation-model versions and inference
settings are recorded; no parameter training occurs in the main
experiment.

The adaptive state consists of an experience collection \(E_t\) and
relationship records \(R_t\). Immediate state \(Z_t\) includes task
progress, outstanding claims, current load, availability, and
permissions. Separating these states is essential: disabling learning
must preserve operational permission checks and changes in availability.

\subsection{Conditions observable before a
decision}\label{conditions-observable-before-a-decision}

Propose a fixed, versioned function \(x_d=F_v(I_d)\) of the member's
authorized pre-decision information \(I_d\). The record retains the
source task/graph/environment versions, extraction position, observation
order, and component-wise missingness. Final patches, later failures,
and sealed evaluation never fill in an earlier \(x_d\). An approved
task-metadata revision affects subsequent decisions and keeps its own
source version; an earlier decision and its relationship key remain
immutable.

The following is one concrete condition catalogue for software-task
development. It is a method proposal, not an approved product schema or
registered set of categories. The same approved catalogue and extraction
code would be fixed across all eight configurations.

{\def\LTcaptype{none} % do not increment counter
\begin{longtable}[]{@{}
  >{\raggedright\arraybackslash}p{(\linewidth - 4\tabcolsep) * \real{0.3333}}
  >{\raggedright\arraybackslash}p{(\linewidth - 4\tabcolsep) * \real{0.3333}}
  >{\raggedright\arraybackslash}p{(\linewidth - 4\tabcolsep) * \real{0.3333}}@{}}
\toprule\noalign{}
\begin{minipage}[b]{\linewidth}\raggedright
Component
\end{minipage} & \begin{minipage}[b]{\linewidth}\raggedright
Proposed extraction from pre-decision sources
\end{minipage} & \begin{minipage}[b]{\linewidth}\raggedright
Proposed grouping value
\end{minipage} \\
\midrule\noalign{}
\endhead
\bottomrule\noalign{}
\endlastfoot
Dependency structure & Count direct predecessors in the approved
task-graph version; retain the full declared predecessor set & Zero,
one, or multiple predecessors \\
Potential write overlap & Apply the project's existing path
normalization and overlap rule to declared write scopes in that graph,
before inspecting actual patches & Declared overlap present or absent \\
Required tools & Read the explicitly declared tool/capability families
from the task version; retain their version requirements & Exact finite
set of families in the approved catalogue \\
Modification kind & Read a declared task label, rather than classify the
completed change & Documentation, configuration, source, tests, mixed,
or other declared kind \\
Environment & Read the declared language/runtime family and required
environment version & Exact approved environment class \\
Verification obligation & Read the permitted feedback-check or
human-rubric kind and its version, without sealed answers & Exact
approved obligation class \\
\end{longtable}
}

Every component also has a distinguished unavailable value \(\bot\). An
explicitly declared empty tool set differs from an unavailable
declaration. Missing paths do not mean disjoint paths, and unknown
compatibility does not mean compatibility. A malformed value or a known
value outside the approved catalogue stops extraction for the affected
decision. A permitted missing value remains \(\bot\); it is not imputed
from outcomes or a selected collaborator's performance. Missing
information required for authorization, capability, or safe execution
prevents assignment under the operational gate.

Define the condition scope \(c_d=g_v(x_d)\) as the exact tuple of these
grouping values together with the approved project/sharing scope and
extractor version. Distinct missingness patterns remain distinct groups.
This rule uses no learned outcome-dependent boundaries and no automatic
pooling into a neighboring group. Scope identifiers denote this semantic
tuple; this proposal does not require a new hashing or storage protocol.
A different catalogue is a new method version, not an unrecorded
within-study adjustment.

The versioned \(x_d\) supports both experience-applicability checks and
relationship grouping. Any additional structural or semantic retrieval
score uses its declared frozen representation and is recorded
separately. Its model, normalization, and utility coefficients remain
registration fields. Current load, active write conflicts, availability,
and live permissions stay in \(Z_d\) and are checked for each decision;
the historical condition group cannot replace them. The proposed
catalogue and grouping have not been implemented or empirically
validated in the inspected fixed policy.

\section{Proposed method}\label{proposed-method}

\subsection{Versioned experience and use
records}\label{versioned-experience-and-use-records}

An experience unit records the task conditions, suggested actions,
collaboration arrangements, participant capability versions, source
events, verification evidence, failure conditions, cost, use history,
and version lineage. Task experience describes problem-solving or
tool-use procedures. Collaboration experience describes role
assignments, handoff requirements, and checks on contributions.

A retrieval hit establishes availability of an experience. Actual
adoption requires a use record linking its exact version to a submitted
decision and subsequent execution. This distinction permits measurement
of exposure, adoption, and downstream outcomes without counting all
retrieved entries as used.

The research draft distinguishes unverified, single-task-verified, and
multi-condition-verified experience. Unverified entries enter an
isolated verification queue. Ordinary decisions retrieve only entries
meeting the applicable evidence requirement \textbf{and} current
authorization. In the development defaults, one independent successful
task supports single-task verification; at least three independent tasks
covering two predefined conditions support multi-condition verification.
These are draft operating rules, with no claim that three tests
establish general reliability.

An applicable counterexample suspends ordinary recommendation and
prompts revision and revalidation. Repeated checks on one task do not
create independent task coverage. Evidence records distinguish automated
tests, checkable tool results, blinded human evaluation, and model
judgments.

Evidence strength and permission to use an entry are separate concepts.
Existing human approval must remain effective. A future evidence-level
implementation cannot automatically grant the current product's approval
or cross-project sharing permission. This separation requires an
explicit product decision before execution.

\subsection{Directed relationships and specific
obligations}\label{directed-relationships-and-specific-obligations}

Let a persistent relationship key be \(\rho=(i,v_i,j,v_j,h,o,c)\):
requester and collaborator identities, their known capability versions,
collaboration type \(h\), obligation class \(o\), and condition scope
\(c\) captured at the originating decision. Obligation classes permit
accumulation over comparable tasks. A specific obligation instance
additionally binds the actual task, handoff, target artifact/source
version, required checks, and resolution rule. It is not interchangeable
with its class.

Use one outcome unit \(u\) per task and relationship key. Fix its
aggregation and closure rule at first use; each constituent handoff
declares its target obligation before execution. A terminal pass
requires the declared target obligations to pass; a terminal failure
requires an independently checkable breach under that resolution rule.
An unresolved check, unclosed retry, or uncertain attribution leaves the
unit unadjudicated. Several tests or retries inside the same task/unit
do not create several complete relationship counts. Their executions and
costs remain recorded. Different obligation classes require separate
predeclared keys, not classes invented after observing which check
failed.

While task work for a unit remains open, each real retry or new target
retains its own source and artifact version under the same declared
resolution rule. At the task-resolution cutoff, that rule fixes which
final target versions the single unit verdict assesses; an unfinished
independent check can leave adjudication pending after this cutoff.
Subsequent corrections reassess those same fixed targets and sources.
New work or a new artifact after closure is not a correction that
overwrites the earlier target, and cannot manufacture an additional
count for the same task/key. Its treatment requires the predeclared
task/obligation rules, with every execution retained.

A verdict links the exact obligation instance, its captured \(\rho\),
target version, actual source event, check/rubric version, result
\(r_u\in\{0,1\}\), and evidence weight \(q_u\in[0,1]\). The development
primary rule uses \(q_u=1\) for an admissible independent target check.
Model opinion, executor self-report, ordinary process exit, and
unresolved attribution supply no successful or failed count. Human
checks identify the target and use a fixed configuration-blinded rubric.
Task-level success does not credit every relationship in a chain, and an
infrastructure interruption is not a breach by every collaborator.
Attribution to a particular obligation identifies what was checked; it
does not identify that collaborator's causal contribution to total task
success.

These are semantic research definitions to be mapped to existing task,
operation, artifact, and event identities. This manuscript creates no
events or replacement execution authority. Evidence strength, access to
the source, and current human permission to reuse experience remain
separate requirements.

\subsection{Scheduled position, visibility, and
corrections}\label{scheduled-position-visibility-and-corrections}

Use the predeclared task position \(t=1,\ldots,T\) as the decay clock.
It advances with the common scheduled sequence, including failed or
unattempted scheduled positions, rather than the number of successful
completions, model calls, messages, or relationship updates. Wall-clock
time and waiting are measured separately. The index uses existing task
order and does not introduce a scheduler or timer service.

For an assessed outcome, \(b_u\) is the logical position of the
underlying obligation outcome/target version. For verdict record \(e\),
\(a_e\) is the first logical position at which that record and its
required provenance are validated and available at the existing
authoritative record source for the authorized update process. This is
not the time at which every member has read it. Require \(b_u\leq a_e\).
Missing or contradictory timing/provenance stops the affected update; an
arrival time is not substituted for an unknown \(b_u\). A correction of
the same target retains \(b_u\). A new target is not silently renamed as
a correction of an older one.

A decision at position \(t\) uses only verdicts with \(a_e<t\). Newly
available verdicts at \(t\) wait until a later position, including when
execution within \(t\) has several decisions. This is a necessary
visibility gate, not a guarantee that the decision has received all
eligible records. Record the existing authoritative within-position
order for arrivals and decisions. For records first available at the
same position, the authoritative order determines the correction chain;
transport arrival order is not a replacement for it. Live availability,
permissions, and claim state may change immediately within the position
through \(Z_d\).

For each unit, an initial adjudication and explicitly linked authorized
corrections form one validated chain. Unverified opinions or pending
check reports do not enter that chain and cannot replace an existing
verified verdict. A withdrawal identifies the verdict, source and
authorized reason it invalidates. Duplicate deliveries retain the same
evidence identity and first-availability position. Repeated test reports
for a unit remain evidence for its single adjudication.

Let \(V_d\) describe the version and actual record collection read and
authorized for decision \(d\) before submission. It preserves existing
source identities/versions, snapshot production position or query scope,
the authoritative order and parent links represented, and the actual
observation/read time. A version label alone is insufficient: the
retained source material must identify the collection and permit
checking its contents. This notation describes an input view, not a new
version service. Let \(H_d\) be its admissible verdict records:
source/provenance and current access are valid, \(a_e<t\), and each
included correction's required parents are present in that view or its
provenance-preserving initial snapshot. A missing parent, known omitted
eligible successor with \(a_e<t\), inconsistent target, contradictory
time, or competing incomparable corrections stops the affected
score/assignment. The process cannot silently choose an older convenient
chain.

For an enabled relationship channel, select each unit's unique latest
verdict within this actual parent-closed collection and authoritative
order. A latest withdrawal contributes nothing; a unit with no available
adjudication remains pending. A known eligible withdrawal/correction
cannot be omitted to retain an older score: incorporate its validated
closed chain or stop the affected assignment. A later correction not
actually read in \(V_d\) is unknown to this decision; record the
observed source version and its real age without claiming a globally
latest state. Later read-only audit may establish additional lag, but it
cannot supply future knowledge to the original decision. Current
permission revocation is checked independently even when a learned
channel is frozen.

Let \(\mathcal U_\rho(t;V_d)\) denote active units for \(\rho\) selected
from \(H_d\). Counts below use this actual set and the current decision
position \(t\), not the clock position at which a cached snapshot was
produced. A stale but complete permitted snapshot can omit unknown later
evidence; it still ages its represented units from their original
\(b_u\) to the present \(t\) when relationship updates are enabled. An
aggregate without the required unit decomposition cannot establish the
specified correction/age semantics. The disabled channel instead uses
its position-zero projection, with no such re-aging. Late verdicts and
corrections alter later views; they never rewrite a submitted decision,
its captured conditions, or its reported propensity. If an old condition
record is invalidated, withdraw its contribution prospectively under the
enabled projection; do not relabel it into a group chosen using later
outcomes.

\subsection{Discounted state and
initialization}\label{discounted-state-and-initialization}

For \(0<\lambda\leq1\), define nonnegative counts from the effective
visible verdicts:

\[
s_\rho(t;V_d)=\sum_{u\in\mathcal U_\rho(t;V_d)}
q_u r_u\lambda^{t-b_u},\qquad
f_\rho(t;V_d)=\sum_{u\in\mathcal U_\rho(t;V_d)}
q_u(1-r_u)\lambda^{t-b_u}.
\]

With fixed \(\alpha_0,\beta_0>0\), set

\[
w_\rho(t)=\frac{\alpha_0+s_\rho(t)}
{\alpha_0+\beta_0+s_\rho(t)+f_\rho(t)}.
\]

Here and below, \(s_\rho(t),f_\rho(t),w_\rho(t)\) suppress \(V_d\) when
its identity is fixed explicitly by the decision. Equal \(t\) alone does
not imply equal views or counts.

Each unit appears once with its effective verdict. Replacing or
withdrawing a verdict recomputes its contribution to this finite sum; it
does not append an extra success/failure. Delayed evidence is discounted
from \(b_u\), not rejuvenated at \(a_e\). With the same represented
effective units and verdicts and no new evidence, both counts multiply
by \(\lambda^{t'-t}\) between \(t\) and \(t'\). The original one-event
recursion is recovered by adding a newly visible unit's already-aged
term to those decayed counts. A zero-weight record adds no evidence
while elapsed-position decay can still occur.

At position zero, the common nonempty historical snapshot contains
approved experience versions and historical units with source, target,
effective verdict, and documented occurrence positions
\(b_u\leq a_e<0\). Pre-study history uses the same registered ordinal
clock with zero as the study boundary; its verdicts and provenance are
available before that boundary. Snapshot counts are the same sums
evaluated at zero. Restoring a snapshot retains unit identities and
correction lineage, so historical events are not counted again as new
evidence. Initial nonzero aggregates require per-unit
provenance/decomposition for later withdrawal or correction;
aggregate-only history is insufficient for this proposal. Historical
positions cannot be invented to retrofit an existing store. An enabled
view retaining the same initial units and unchanged verdicts gives their
contribution at \(t\) as exactly \(\lambda^t s_\rho(0)\) and
\(\lambda^t f_\rho(0)\).

The development illustrations retain \(\alpha_0=\beta_0=1\) and
\(\lambda=2^{-1/20}\), interpreted as a half-life of 20 scheduled
positions. These are unregistered values. A newly eligible, fully
identified capability version with no evidence has zero counts and the
defined prior-like score \(\alpha_0/(\alpha_0+\beta_0)\). Propose no
automatic transfer across capability versions or condition-schema
versions. An unannounced change remains unobserved until evidence
reveals it; earlier decisions keep the capability information they
actually had. No version or evidence update changes current permissions.

When relationship updates are off, decisions read the position-zero
counts and weight without decay, new verdicts, or corrections entering
the learned state. Logs remain writable. A live withdrawal of permission
still blocks unauthorized reads/use through the operational gate;
freezing a learned representation cannot authorize it. When experience
updates are off, content, evidence level, historical use/ranking inputs,
and associated adaptive caches remain frozen. Enabled experience
revisions similarly become eligible only after their declared
availability position and current approval checks.

These scores summarize selected, discounted evidence. Unique task
accounting does not make observations independent, and the prior-like
ratio is not a calibrated Bayesian posterior. The complete-block
analysis and its sampling limitations remain those in Section 6.

\subsection{Common decision rule}\label{common-decision-rule}

For eligible candidates \(C\), the proposed \texttt{policy\_score}
combines normalized task fit, experience support, relationship score,
estimated cost, and current load. Development defaults add the first
three and subtract the last two, with coefficients one. Accurate feature
definitions, aggregation, and cost prediction remain registration
fields.

For candidate \(j\) with score \(U_j\), choose according to

\[
p(j)=(1-\epsilon)
\frac{\exp(U_j/\tau)}{\sum_{k\in C}\exp(U_k/\tau)}
+\frac{\epsilon}{|C|}.
\]

The draft defaults are \(\epsilon=0.10\) and \(\tau=1\). Candidate
ordering, score components, probabilities, random draws, selected
experience versions, observation time, state version, submission time,
and claim outcome must be observable. An empty candidate set produces an
unassigned decision under the common waiting rule. Exploration operates
inside capability and authorization constraints.

Required but undefined condition, permission or utility inputs prevent
the affected assignment with an explicit error. They are not replaced
with guessed zero scores or an unrecorded alternative policy. The
defined zero-evidence state for a known, eligible relationship key is
separate from missing identity or unreadable mandatory state. Current
authorization is checked again by the common claim/execution service at
submission.

Central and local decisions call the same implementation. For a decision
assigned to member \(i\), the central service receives that member's
actual information/record view \(I_d,V_d\), candidate set, retrieval
allowance, and decision budget; access to an authoritative source does
not add records to that designated input. Both use the same claim and
conflict-handling service and the same total decision-concurrency
allowance. The central service submits choices for members; local
members submit their own choices.

\subsection{Conditional properties and
proofs}\label{conditional-properties-and-proofs}

The following properties describe the proposed definitions under stated
assumptions. They are elementary algebra and coupling arguments, not a
claim of original general theory, empirical benefit, calibrated
prediction, regret control, or causal identification. The inspected
fixed policy has not acquired these properties merely because the
specification states them.

\subsubsection{Canonical evidence and duplicate
delivery}\label{canonical-evidence-and-duplicate-delivery}

\textbf{Claim.} Fix a decision position, relation key, the actual
admissible collection \(H_d\) from a specified \(V_d\), its
authenticated parent-closed correction chains, and its occurrence
positions. Assume every unit has one unambiguous effective verdict or
withdrawal. The counts are invariant to the order in which those records
are enumerated and to repeated delivery of a record with the same
evidence identity.

\textbf{Proof.} Each unit's effective verdict is a function of the fixed
visible chain and authoritative order, not the enumeration order.
Repeating an identity does not change that chain. Each active unit
therefore supplies the same single two-component contribution. Finite
addition is commutative and associative, so the two sums and their ratio
agree. This compares the same visible set. Different first-availability
positions can change an online decision and are not covered by the
claim.

\subsubsection{Sensitivity of a relationship
score}\label{sensitivity-of-a-relationship-score}

\textbf{Claim.} Compare two nonnegative states for the same relation key
and the same positive priors. Put \(m_0=\alpha_0+\beta_0\). Then

\[
|w-w'|\leq\min\left\{1,
\frac{|s-s'|+|f-f'|}{m_0}\right\}.
\]

The following evidence-mass refinement additionally assumes that both
states are generated by the enabled relationship-update projection at
the same position \(t\), with the same ordinal clock and \(\lambda\).
For each unit in the union of their effective sets, let
\(v_u=(q_u r_u\lambda^{t-b_u},q_u(1-r_u)\lambda^{t-b_u})\) be its actual
contribution under that state's effective verdict, using the zero vector
when absent; define \(v'_u\) from the other state's effective verdict.
Thus \((s,f)=\sum_u v_u\) and \((s',f')=\sum_u v'_u\). If
\(M=\sum_u\|v_u-v'_u\|_1\), the numerator above is at most \(M\). This
refinement does not compare an enabled, decaying state with a
relationship-off frozen snapshot. That comparison can use the original
bound with each projection's actual \(s,f\); identical historical
identities need not then have identical mass.

\textbf{Proof.} Write \(D=m_0+s+f\). The derivatives of
\(w=(\alpha_0+s)/D\) are \((\beta_0+f)/D^2\) and \(-(\alpha_0+s)/D^2\),
whose absolute values are at most \(1/m_0\). The segment joining the two
nonnegative states stays in that domain. Integrating the derivative
along the segment gives the stated Lipschitz bound; \(0<w,w'<1\) gives
the other bound. The triangle inequality applied to the contribution
vectors gives the \(M\) bound. Missing evidence has its own discounted
mass; a success-to-failure correction changes both components and can
contribute twice that mass to \(M\). The bound does not treat such a
correction as a single missing count.

\subsubsection{Sensitivity of the fixed-candidate choice
distribution}\label{sensitivity-of-the-fixed-candidate-choice-distribution}

\textbf{Claim.} Let the same nonempty finite candidate set, ordering,
\(\tau>0\), and \(0\leq\epsilon\leq1\) be used in two decisions. If all
utilities are finite and \(\|U-U'\|_\infty\leq\eta\), their
distributions satisfy

\[
\operatorname{TV}(p,p')\leq
(1-\epsilon)\min\{1,\eta/\tau\}.
\]

\textbf{Proof.} Along \(U(z)=U+z(U'-U)\), let \(\sigma_j(z)\) denote the
softmax term and \(\Delta=U'-U\). Differentiation gives

\[
\frac{d\sigma_j}{dz}=
\frac{\sigma_j}{\tau}\left(\Delta_j-
\sum_k\sigma_k\Delta_k\right).
\]

Because \(|\Delta_j|\leq\eta\), the sum of absolute derivatives is at
most \(2\eta/\tau\). Integrating and dividing by two bounds total
variation by \(\eta/\tau\); total variation is also at most one. The
identical uniform exploration component cancels, leaving the factor
\(1-\epsilon\).

If only an affine relationship term \(c_R w_\rho\) changes and all other
utility inputs/scales are fixed, the preceding evidence-mass refinement
supplies \(\eta\leq |c_R|\max_\rho\min\{1,M_\rho/m_0\}\) only when both
states meet its same-position, same-clock, enabled-projection
assumptions. For an enabled/frozen comparison, use each projection's
actual counts in the original score bound to obtain a utility bound
instead. A changing normalization must be included in the utility error.
Candidate entry/removal, changed authorization, different retrievals, or
unequal budgets violate the fixed-input premises; the bound then
supplies no task-quality or deployment conclusion.

\subsubsection{Equal information and equal random
draws}\label{equal-information-and-equal-random-draws}

\textbf{Claim.} With the same deterministic policy implementation,
eligible ordered candidates, identical actual information/state views
\(I_d,V_d\) and their admissible collections/projections, finite
utilities, and the same uniform draw, central and local decision makers
select the same candidate.

\textbf{Proof.} Their score vectors and probability vectors are
identical. Their cumulative probabilities in the common ordering are
identical, so inverse-CDF selection with a fixed boundary convention
yields the same candidate for the same draw. Induction extends the
statement over decisions only if initial states and all intervening
exogenous inputs, evidence availability/corrections, authorization and
claim outcomes are also coupled identically. Different provider outputs
or request/claim timing are not made equal by a requested seed.
Synchronous replay checks this conditional equivalence; it cannot
establish a positive effect of location or equal realized performance
under a different deployment.

\subsubsection{History excluded by disabled
updates}\label{history-excluded-by-disabled-updates}

\textbf{Claim.} Fix task inputs, operational state, initial snapshots,
policy/random draws, and the enabled channel's allowed state. Changing
only the post-initialization history of a disabled channel does not
change decision probabilities under the specified state projection.

\textbf{Proof.} For a disabled relationship channel, the policy receives
the initial counts/weight with no decay or historical corrections. For a
disabled experience channel, it receives initial content,
verification/use/ranking values and adaptive cache inputs. The changed
later history enters neither projection. All remaining policy arguments
are fixed, so the probability vector is fixed. This is conditional on
enforcing the projection at every read, including caches and summaries.
If another enabled channel actually changes, or current
permissions/availability change, the premises do not hold. This
specification is not evidence that the product already enforces those
exclusions.

\subsection{Feedback cycle and failure
handling}\label{feedback-cycle-and-failure-handling}

After an executed task, the allowed feedback suite produces
obligation-specific evidence. The experience process proposes and
verifies revisions within its assigned budget. Validated verdicts,
withdrawals and corrections enter their recorded availability positions.
Later decisions use the effective earlier-visible set and the projection
allowed by their configuration. The intended cycle is specified below;
it is not an executed algorithm trace.

\begin{Shaded}
\begin{Highlighting}[]
\NormalTok{For each scheduled task in an isolated configuration:}
\NormalTok{    obtain current operational state and permitted earlier{-}visible adaptive state}
\NormalTok{    capture versioned pre{-}decision conditions, obligations, and eligible candidates}
\NormalTok{    retrieve eligible experience and compute the common decision probabilities}
\NormalTok{    submit the selected task or partner decision through the common claim service}
\NormalTok{    execute within the registered task and sequence budgets}
\NormalTok{    record actual feedback sources, occurrence/availability, usage, and unresolved units}
\NormalTok{    verify experience and adjudications under current permissions and declared budgets}
\NormalTok{    if updates are enabled: expose valid revisions at their allowed future positions}
\NormalTok{    submit the final artifact to the isolated evaluator}
\end{Highlighting}
\end{Shaded}

Unavailable shared services prevent affected new decisions and produce
an explicit failure or waiting record. Recovery must follow the
registered policy, including any explicitly selected snapshot mode and
its version. No unobserved retry, substituted state, or unknown-delivery
resend can be treated as successful execution. Duplicate evidence
handling reuses event identity; it does not create a separate scheduling
authority.

\section{Experimental protocol draft}\label{experimental-protocol-draft}

\subsection{Eight paired
configurations}\label{eight-paired-configurations}

Every replication contains the following eight configurations in
randomized execution order. All begin from the same nonempty experience
and relationship snapshots, copied into isolated condition state, and
process the same scheduled task sequence.

{\def\LTcaptype{none} % do not increment counter
\begin{longtable}[]{@{}
  >{\raggedright\arraybackslash}p{(\linewidth - 6\tabcolsep) * \real{0.2500}}
  >{\raggedright\arraybackslash}p{(\linewidth - 6\tabcolsep) * \real{0.2500}}
  >{\raggedright\arraybackslash}p{(\linewidth - 6\tabcolsep) * \real{0.2500}}
  >{\raggedright\arraybackslash}p{(\linewidth - 6\tabcolsep) * \real{0.2500}}@{}}
\toprule\noalign{}
\begin{minipage}[b]{\linewidth}\raggedright
Configuration
\end{minipage} & \begin{minipage}[b]{\linewidth}\raggedright
Decision location
\end{minipage} & \begin{minipage}[b]{\linewidth}\raggedright
Experience updates
\end{minipage} & \begin{minipage}[b]{\linewidth}\raggedright
Relationship updates
\end{minipage} \\
\midrule\noalign{}
\endhead
\bottomrule\noalign{}
\endlastfoot
L00 & Local & Off & Off \\
L10 & Local & On & Off \\
L01 & Local & Off & On \\
L11 & Local & On & On \\
C00 & Central & Off & Off \\
C10 & Central & On & Off \\
C01 & Central & Off & On \\
C11 & Central & On & On \\
\end{longtable}
}

All conditions can retrieve initial experience. Experience-off freezes
content, verification state, historical use statistics, and ranking
terms derived from those statistics. Relevance to a current task may be
recomputed from the frozen representation. Relationship-off freezes
counts, weights, and time decay. Immediate task state, permissions, and
availability continue to change in every condition. Logs remain
writable, but a disabled channel cannot update its decision inputs from
those logs through caches, indexes, summaries, or shared services. An
enabled experience channel can still learn from permitted feedback while
relationship counts are frozen; an enabled relationship channel can
still update while the experience representation is frozen.

Each isolated copy includes historical unit identities, effective
verdicts, occurrence positions, capability/extractor versions, and
correction lineage, not only aggregate counts. New visible
verdicts/revisions are gated separately by the two switches. The
disabled-channel property above is conditional on holding the other
policy inputs fixed; the eight configurations do not assume that an
enabled channel or operational trajectory remains unchanged. Current
permission withdrawal blocks use without silently rewriting frozen
historical weights.

Model lists, prompts, source collections, feature extraction, task order
within each paired block, concurrency quotas, and task/sequence budget
rules are shared. New literature or other knowledge cannot enter only
selected conditions. Initialization and execution order are recorded
together with service versions and calendar time to expose drift and
shared-service interference.

\subsection{Decision location and
identification}\label{decision-location-and-identification}

A synchronous development replay supplies equal information snapshots,
event order, candidate identities, and random draws to both
implementations. Matching selections test implementation consistency.
With equal state and random draws, the common rule implies equal
decisions; replay supplies no evidence for a positive location effect.

The deployed comparison records each actual \(V_d\), its represented
source collection, snapshot production position, read/submission times,
request waiting, claim ordering/conflicts, and communication. Observed
snapshot age and authoritative first availability are distinct; unknown
later corrections cannot be labeled known at the decision. Different
actual views can yield different scores even at the same \(t\), while an
enabled projection always ages its represented evidence to that
decision's \(t\). A nonzero location effect would describe those
specified deployment conditions. It cannot be attributed to abstract
decentralization while central caching, extra information, different
scoring, or unequal compute also vary. Shared storage, process
placement, cache policy, communication path, and scheduling must
therefore be declared. Whether the deployment introduces a meaningful
location intervention remains a readiness question.

\subsection{Tasks, feedback, and sealed
evaluation}\label{tasks-feedback-and-sealed-evaluation}

Each task has a \texttt{feedback\_suite} and an
\texttt{evaluation\_suite}. Feedback is available to execution and the
permitted updates. Propose keeping formal evaluation programs, answers,
item-level outcomes and scores sealed from those processes until
\textbf{all planned paired blocks of the main study} have completed.
Finishing one sequence or one block does not release its evaluation.
This stronger release boundary prevents early blocks from influencing
later execution, experience, task selection, or analysis settings. It
remains a proposal requiring approval and actual enforcement.

\subsubsection{Roles, fixed artifacts, and blinded
evaluation}\label{roles-fixed-artifacts-and-blinded-evaluation}

The proposed access matrix below describes information permissions
before that release. Execution includes participating agents and their
coordinators; updates include experience distillation, retrieval/ranking
maintenance, and relationship adjudication. The independent evaluator
includes any human raters. The analyst has no execution or update role.
Permissions refer to study material, not permission to inspect unrelated
user data.

{\def\LTcaptype{none} % do not increment counter
\begin{longtable}[]{@{}
  >{\raggedright\arraybackslash}p{(\linewidth - 8\tabcolsep) * \real{0.2000}}
  >{\raggedright\arraybackslash}p{(\linewidth - 8\tabcolsep) * \real{0.2000}}
  >{\raggedright\arraybackslash}p{(\linewidth - 8\tabcolsep) * \real{0.2000}}
  >{\raggedright\arraybackslash}p{(\linewidth - 8\tabcolsep) * \real{0.2000}}
  >{\raggedright\arraybackslash}p{(\linewidth - 8\tabcolsep) * \real{0.2000}}@{}}
\toprule\noalign{}
\begin{minipage}[b]{\linewidth}\raggedright
Material
\end{minipage} & \begin{minipage}[b]{\linewidth}\raggedright
Execution
\end{minipage} & \begin{minipage}[b]{\linewidth}\raggedright
Updates
\end{minipage} & \begin{minipage}[b]{\linewidth}\raggedright
Evaluator
\end{minipage} & \begin{minipage}[b]{\linewidth}\raggedright
Analyst before release
\end{minipage} \\
\midrule\noalign{}
\endhead
\bottomrule\noalign{}
\endlastfoot
Approved task inputs and feedback checks/results & Assigned scope &
Authorized feedback & Task/rubric context & Sanitized progress only \\
Final fixed artifact and required environment & Submit own version &
Feedback-authorized sources only & Read submitted version & Submission
identity only \\
Sealed evaluation program and target answers & No & No & Evaluate only &
No \\
Sealed item outcomes, scores and evaluation traces & No & No & Retain
privately & No \\
Configuration mapping and decision histories & Assigned execution scope
& Assigned switches/sources & Hidden & Mapping withheld \\
Model/tool usage, costs and operational errors & Own authorized
operations & Own authorized operations & Own evaluation records &
Sanitized accounting \\
\end{longtable}
}

The evaluator receives the final immutable artifact/source version, task
and permitted rubric context, and an evaluation label whose
configuration mapping is retained separately. Its presentation removes
irrelevant configuration labels and execution commentary while
preserving the unaltered artifact and necessary execution context.
Source/version linkage remains auditable. Human raters use a fixed
rubric without update switches, decision location, partner histories,
preliminary scores, or other raters' judgments. The design cannot
guarantee blindness when the artifact itself reveals the condition;
record unavoidable cues and any unblinding. Assigned execution agents
may know their own switches, so this is evaluator blinding, not a claim
that every role is blinded.

Submission acknowledgment confirms receipt of the fixed version and
returns no grade, assertion failure, evaluation exit status,
result-dependent timing report, or diagnostic content. Evaluation writes
neither task feedback nor shared learning state. Final artifacts cannot
be revised after their sealed evaluation is inspected. Infrastructure
errors and incomplete evaluation remain explicit in the evaluator's
compartment; outcome-bearing messages cannot be used as ordinary
operational feedback.

Before formal collection, propose fixing the task/suite versions, full
block count, analysis procedure, damaged-task rules and release rule. At
the release boundary, all actually produced final artifacts are fixed
and submitted, and every planned position has its sourced terminal
disposition. A budget-exhausted or otherwise observed algorithmically
unfinished position may have no artifact: retain that absence and its
actual reason, with the failure outcome assigned by the predeclared
protocol rule. Do not manufacture an artifact or claim a sealed-suite
measurement for that position. These sourced algorithmic failures differ
from an unknown outcome caused by missing evaluation, damaged data or an
incomplete block; the latter remain unavailable and stop the affected
paired analysis in Section 6. Unknown measurements are not filled with
zero to release results. Accounting unknowns remain labeled, and
exclusion decisions do not use sealed scores. After all planned blocks
have completed their required collection and dispositions, the analyst
receives the frozen outcome records, configuration mapping and
accounting sources. An unfinished study does not automatically release
partial scores or authorize outcome-based replacement blocks; its
closure requires an explicit documented decision.

\subsubsection{Leakage paths and required access
checks}\label{leakage-paths-and-required-access-checks}

The sealed compartment must exclude execution and update reads through
files, repositories/attachments, database queries, HTTP responses,
process output, ordinary logs, retrieval indexes, summaries, caches and
model context. Error text, exit codes, test names, timing details and
notifications can also reveal evaluation outcomes. An acknowledgment
path must remain independent of the score. Evaluators must not share a
writable index, conversation or learned cache with participants. The
same restrictions apply to human coordination messages and automated
experience proposals.

Role names and separate directories alone provide no access-isolation
evidence. A future authorized implementation must verify actual read
denial and result-free responses across each relevant route using real
fixed public-project artifacts and actual source events; each claimed
event must be locatable in the real event store. It must also check that
repeated submissions, errors, retrieval and restored state do not reveal
a sealed item. No such tests, events or permissions are created in this
writing round. The current native file/command-tool restriction and
unresolved authority choices remain in force; this matrix does not
select a product component or grant a new role.

Task and suite selection are unfinished. The registration record must
identify exact versions, lawful source access, suite overlap, and the
public-project scope. A content audit must check whether public feedback
already exposes the sealed assertions or target answers. Access
isolation cannot remedy that duplication. Independent pilot tasks and
their development evaluation are separated from confirmatory collection;
revealing pilot results for planning does not authorize reuse of formal
sealed results.

The primary task outcome is binary satisfaction of the fixed evaluation
suite. Errors, timeouts, budget exhaustion, and unfinished tasks count
as failures. After a sequence exhausts its budget, all remaining
scheduled tasks remain in the denominator. Damaged task data or missing
evaluation programs follow predeclared, configuration-symmetric
exclusion rules. Algorithmic failures, unhelpful experience, invalid
tools, waiting, and retries remain outcomes and costs. A retry never
erases its original execution record.

\subsection{Budgets and accounting}\label{budgets-and-accounting}

Accounting covers model calls, tools, experience extraction, retrieval,
verification, coordination, retries, initialization, and independent
evaluation. Record these categories through existing task/operation
identities and actual controller, provider, tool and payment sources.
This proposes a research interpretation of those records, not a new
Attempt entity, scheduler or accounting service. The
unverified-experience queue retains its draft ceiling of 5\% of a
configuration's total budget, with its work charged to that
configuration. Neither that number nor a formal budget unit is approved.

\subsubsection{Observed operations, retries, and inclusion
relationships}\label{observed-operations-retries-and-inclusion-relationships}

A scheduled task position is an outcome-denominator item. It may have no
submitted request after budget exhaustion, or several genuine operations
and retries. A platform submission, controller turn and
provider-observable attempt are different observations; retain their
actual links when the sources expose them. A platform request does not
establish that the provider made exactly one internal attempt.
Unobservable internal usage remains unknown.

For each cost observation, retain the original operation/source identity
and version, task and configuration linkage, model/tool/service version,
observation time, raw quantity and unit, and whether it is an inclusive
total or a component. Repeated transport of one usage record is one
observation. A cumulative or corrected report for the same operation
replaces the earlier reported total for accounting; it is not an extra
call to add. Preserve the earlier source and correction history.
Distinct actual billable attempts remain distinct costs, including
failed, timed-out and unhelpful attempts.

Nested-agent reports and coordinator summaries can already contain child
operations. Establish the documented inclusion relation before summing:
count a disjoint set of actual source quantities once, retaining rollups
as cross-checks. Do not add a parent's inclusive total to its children,
or sum an increment and the cumulative total containing it. Unknown
overlap leaves that aggregate unresolved. The same rule applies to usage
components: cached input or reasoning counts are added only when the
actual interface defines them as disjoint from the reported input/output
totals. Their names alone do not establish disjointness.

Unknown delivery or missing usage never triggers a resend for
accounting. Reconcile against the original controller/provider records
and later legitimate receipts without reissuing the original operation.
Missing receipts, failed execution and subscription coverage do not
establish zero consumption. An actual zero needs source evidence, such
as an unsubmitted operation, rather than absence of a record. Report
known quantities and their coverage separately from unresolved
quantities. The unresolved S3 call-count unit and tool/destination
evidence choices still prevent a formal budget-matched claim for the
affected path; this proposal selects no default.

\subsubsection{Payments, resource quantities, and elapsed
time}\label{payments-resource-quantities-and-elapsed-time}

Model usage reports input, cached input, output and reasoning quantities
only with their actual inclusion semantics, exact model version and
observation/billing period. Tools, retrieval, distillation, verification
and coordination retain their observed units separately. Do not add
tokens across incompatible models and call that a common resource budget
without an approved conversion. Rates and payment attribution require
actual applicable billing evidence; subscription payments and allocation
rules remain distinct from per-request resource consumption. Unknown
payment is not zero payment.

Record communication messages and bytes, and model-context tokens when
actually metered; a token estimate must identify its estimator and
remain labeled as an estimate. Report sequence makespan, waiting,
individual operation intervals and their overlap separately. Summed
concurrent durations are not elapsed study time, and waiting time is not
measured CPU/GPU time. Actual CPU/GPU quantities require their own
meters. Evaluation resources remain separately identifiable and
outcome-bearing evaluator traces remain sealed; sanitized accounting
does not release scores.

\subsubsection{Shared preparation and per-sequence
restoration}\label{shared-preparation-and-per-sequence-restoration}

Let \(P\) be the actual preparation cost of the one common frozen
history used for the main study, recorded by each compatible
resource/payment unit. It includes the actual acquisition, distillation
and verification used to produce that snapshot, not an invented
reconstruction of unavailable history. Propose allocating \(P\) equally
over the fixed planned \(8r\) main-study sequences. For sequence
\(\ell\), report component-wise

\[
C_\ell^{\mathrm{report}}=
\frac{P}{8r}+C_\ell^{\mathrm{restore}}
+C_\ell^{\mathrm{execution}}+C_\ell^{\mathrm{evaluation}}.
\]

The execution term contains its disjoint model/tool, learning,
verification, retrieval, coordination and retry costs; restoration
contains that sequence's actual copying/loading/isolation work.
Evaluation is shown separately from execution. The reporting identity
does not choose which unit or component a formal budget caps. Report raw
\(P\), its proposed share, and each execution/restoration/evaluation
component, with the original source inventory. Charge \(P\) only once in
the study total. Pilot, perturbation and iteration-case costs remain
separate; reused preparation is identified rather than charged again as
a new acquisition.

The allocation denominator is the registered planned sequence count, not
a smaller successful/completed count chosen later. If collection ends
early, report planned versus completed sequences and the remaining
allocation explicitly, while the actual study total still contains the
full \(P\) once. Additional real restoration attempts and failed
preparation work remain observed costs. If historical \(P\) is
unavailable, report that fact and any sourced partial quantities without
manufacturing a total or calling history free. The corresponding
end-to-end and budget-matched comparison remains unavailable. Formal
\(r\), cost units, allocation approval and caps all remain undecided.

Let the main study use \(r\) replications of \(T\) tasks, perturbations
use \(r_p\) replications of \(T_p\) tasks in three configurations and
four scenarios, and the case use \(r_s\) replications of three processes
over \(G\) rounds and \(T_s\) evaluations per round. The draft
task-accounting identity is

\[
N_{\mathrm{tasks}}=8rT+12r_pT_p+3r_sGT_s+N_{\mathrm{pilot}}+N_{\mathrm{initial}}.
\]

Case candidate generation and development verification are charged
additionally. The input's \(r=5,T=10\) example gives 40 sequences and
400 task executions; its assumed 30 calls per task gives approximately
12,000 calls before preparation and development work. These are budget
illustrations, with no observed call rate or approved formal sample
size.

\section{Estimands and analysis}\label{estimands-and-analysis}

\subsection{Primary and secondary
contrasts}\label{primary-and-secondary-contrasts}

For replication \(k\), let

\[
Q_{ab,k}^{P}=\frac{1}{T_k}\sum_{u=1}^{T_k}Y_{ab,k,u}^{P},
\qquad P\in\{\mathrm{local},\mathrm{central}\},\quad a,b\in\{0,1\}.
\]

\(T_k\) is the common post-exclusion scheduled denominator across the
eight configurations. We compute each replication's contrasts first,
then average contrasts equally across replications. The planned
position-specific interaction is

\[
I_{\mathrm{feedback}}^{P}=Q_{11}^{P}-Q_{10}^{P}-Q_{01}^{P}+Q_{00}^{P}.
\]

The primary target is \(I_{\mathrm{feedback}}^{\mathrm{local}}\). The
draft practical target \(\delta_{\mathrm{feedback}}=0.03\) means three
additional successes per hundred scheduled tasks beyond the additive
expectation from the individual updates. It does not mean that L11's
total gain over L00 is three points.

Decision-location modification of the interaction is

\[
I_{\mathrm{three}}=I_{\mathrm{feedback}}^{\mathrm{local}}
-I_{\mathrm{feedback}}^{\mathrm{central}}.
\]

Decision-location modification of the total joint-update gain is

\[
I_{\mathrm{position}}=(Q_{11}^{\mathrm{local}}-Q_{00}^{\mathrm{local}})
-(Q_{11}^{\mathrm{central}}-Q_{00}^{\mathrm{central}}).
\]

These are distinct estimands. In each position,

\[
Q_{11}^{P}-Q_{00}^{P}=I_{\mathrm{feedback}}^{P}
+(Q_{10}^{P}-Q_{00}^{P})+(Q_{01}^{P}-Q_{00}^{P}).
\]

Consequently, the location difference in total gain also contains
location differences in the two individual update gains. A large
L11--L00 improvement alone does not establish interaction. We
additionally report L11--L10, L11--L01, and L11--C11 together with
costs. Positive interaction and acceptable operational value are
separate requirements.

\subsection{Replication and
uncertainty}\label{replication-and-uncertainty}

The independent unit is a complete paired block of eight task sequences,
with separately executed requests and reinitialized condition state. For
block \(k\), store the full vector
\(V_k=(Q^{L}_{00,k},Q^{L}_{10,k},Q^{L}_{01,k},Q^{L}_{11,k},Q^{C}_{00,k},Q^{C}_{10,k},Q^{C}_{01,k},Q^{C}_{11,k})\).
Each estimand above is a fixed linear contrast \(D_k=h(V_k)\), estimated
by \(\hat\theta=r^{-1}\sum_{k=1}^{r}D_k\). All blocks receive equal
weight. A pooled task mean would change this estimand when
post-exclusion denominators differ.

Tasks within a sequence are dependent through adaptation. Tests, seeds,
eight configuration labels, and repeated analysis do not add independent
blocks. The nonparametric proposal treats the \(V_k\) as independent,
identically distributed observations from the declared replication
population. Repeated use of a fixed task list and one development
snapshot supports conditional inference for those inputs. Restoring
separate state copies does not independently sample historical
experience. Generalization across tasks or histories requires
corresponding sampling. Shared-provider drift, overlapping requests, or
cross-block state leakage can invalidate the independent-block
assumption; a seed label does not repair it.

Only predeclared, configuration-symmetric damaged-task exclusions alter
a block's common denominator. Budget exhaustion and algorithmic failures
remain scheduled outcomes. Missing measurements or an incomplete block
stop the specified paired analysis; they are not silently deleted,
completed from other conditions, or imputed as observed failures. All
original records remain available for accounting and reporting.

\subsection{Proposed complete-block confidence
procedure}\label{proposed-complete-block-confidence-procedure}

The following BCa procedure is a research proposal, with its
computational settings and suitability to be assessed before formal
collection. It adapts the bias correction, acceleration, and
confidence-point construction of DiCiccio and Efron (DiCiccio and Efron
1996, secs. 2--3 and 6) to the block contrasts defined here. It has not
been executed on Astrocyte pilot or formal outcomes.

Draw \(B\) index vectors of length \(r\) independently with replacement
from \(\{1,\ldots,r\}\). Apply each index vector to entire \(V_k\) rows
and reuse those indices for every contrast. For each contrast retain
every resampled mean \(\hat\theta_b^*\). Let \(\hat G\) be their
right-continuous empirical distribution, with
\(\hat G^{-1}(u)=\inf\{x:\hat G(x)\geq u\}\). A resample with constant
\(D_k\) values remains a legitimate mean. No task-level or
configuration-wise resampling is performed.

Let \(\Phi\) denote the standard normal distribution function and
\(z_u=\Phi^{-1}(u)\). Compute

\[
\hat z_0=\Phi^{-1}\!\left(B^{-1}\sum_{b=1}^{B}
\mathbf 1\{\hat\theta_b^*<\hat\theta\}\right).
\]

For \(r>1\), let \(\hat\theta_{(-k)}\) be the mean omitting complete
block \(k\), let
\(\bar\theta_{(-\cdot)}=r^{-1}\sum_k\hat\theta_{(-k)}\), and set
\(d_k=\bar\theta_{(-\cdot)}-\hat\theta_{(-k)}\). The jackknife
acceleration and adjusted endpoint are

\[
\hat a=\frac{\sum_k d_k^3}{6(\sum_k d_k^2)^{3/2}},
\qquad
A(u)=\Phi\!\left(\hat z_0+
\frac{\hat z_0+z_u}{1-\hat a(\hat z_0+z_u)}\right),
\]

\[
q(u)=\hat G^{-1}(A(u)),\qquad
C_{1-2\gamma}=[q(\gamma),q(1-\gamma)].
\]

The nominal 95\% interval is \(C_{0.95}\) with \(\gamma=0.025\). Its
lower endpoint above the draft \(\delta_{\mathrm{feedback}}=0.03\)
supports the proposed practical primary target; an upper endpoint below
that target indicates it is not reached under the interval rule. An
interval crossing or touching the target is inconclusive. A 95\%
two-sided lower endpoint corresponds to a 97.5\% one-sided bound, which
is the boundary used in the planning example.

The procedure requires a nondegenerate contrast distribution and
applicable independent-block asymptotics. Bounded task scores give
finite moments, but discrete small-block means can violate the smooth
approximations used for higher-order bootstrap claims. No exact
finite-sample coverage or automatic second-order accuracy is asserted
for this design.

Before formal data, specify \(B\), the randomization record, a finite
one-sided significance grid \(\mathcal A\subset(0,1/2]\), and minimum
usable tail ranks. Include the primary level and every one-sided level
needed by the multiplicity decisions below. Require finite \(\hat z_0\)
and \(\hat a\), positive \(\sum_k d_k^2\), positive denominators in
\(A(u)\) for all required \(u\) and \(1-u\), nested finite endpoints on
that branch, and the registered minimum ranks in both tails of every
adjusted quantile. Record distinct block contrasts, ties, extreme
adjusted probabilities, and adjusted quantile ranks. \(B\) and these
numerical criteria remain unconfirmed, not values selected from formal
outcomes.

For computational sensitivity, propose one designated primary Monte
Carlo batch and a fixed number \(J\) of additional independent batches,
each containing \(B\) complete-block resamples. Each batch computes its
own correction and endpoints; within a batch the indices remain common
to all contrasts. Report the endpoint ranges and the grid-inversion
bounds across batches. The primary batch supplies the proposed
inferential quantities; the others diagnose computational instability.
Register \(J\) and endpoint-spread tolerances before formal data. A
diagnostic failure or disagreement about an intended decision makes the
affected inference unavailable; no batch is selected for a favorable
answer and no extra batches are added after inspection. These ranges
describe Monte Carlo sensitivity, not confidence coverage over new
experimental blocks.

Failure of these conditions ends inferential calculation for the
affected contrast. Report the observed blocks and point estimate with
the reason that inference is unavailable. Do not replace a zero
jackknife denominator, clip an infinite correction, remove constant
bootstrap samples, change interval methods, or add formal runs after
seeing the result. A family containing an unavailable marginal
calculation receives no confirmatory family conclusion under this
proposal.

This choice also addresses a specific studentization problem.
Bootstrap-t divides each resampled centered mean by that resample's
standard error (DiCiccio and Efron 1996, sec. 5). If the observed
distinct contrast values have multiplicities \(n_j\), the conditional
probability of a constant resample is \(\sum_j(n_j/r)^r\). Such samples
can have zero standard error even when the original variance is
positive. Deleting, redrawing, or adding a small denominator would
define another procedure. BCa uses their means without that division;
its own degeneracy and tail checks still apply.

\subsection{Null hypotheses and confidence-curve
inversion}\label{null-hypotheses-and-confidence-curve-inversion}

Bootstrap uncertainty around \(\hat\theta\) is not itself a null
distribution. Hall and Wilson (Hall and Wilson 1991, secs. 1--3)
distinguish a resampling comparison centered at the estimate from an
uncentered comparison with the tested value. DiCiccio and Efron
(DiCiccio and Efron 1996, sec. 3) likewise distinguish an ordinary
bootstrap tail area from a hypothesis-testing confidence level. The
proposed tests invert the same adjusted confidence endpoints above. They
do not label a raw fraction of positive bootstrap contrasts as a
calibrated \(p\) value, and do not assume a sharp-null permutation or
exchangeable configuration labels.

For \(H^+_0:\theta\leq\theta_0\), define the grid upper bound
\(\bar p_+(\theta_0)\) as the smallest \(a\in\mathcal A\) such that
\(q(a)>\theta_0\); if none exists, set it to 1. For
\(H^-_0:\theta\geq\theta_0\), define \(\bar p_-(\theta_0)\) using
\(q(1-a)<\theta_0\). The value 1 here records non-rejection over a valid
grid; a failed diagnostic remains unavailable. Strict endpoint
comparisons handle ties without jitter. A value at the smallest grid
level is reported as an upper bound at that resolution, never as an
exact zero \(p\) value.

This is a conservative grid coarsening of inversion of a nested
confidence curve: at a grid level \(a\), \(\bar p_+(\theta_0)\leq a\)
holds exactly when the lower endpoint excludes \(\theta_0\) in the
required direction. The analogous statement holds for the upper
endpoint. Its statistical calibration inherits the approximation of the
BCa confidence curve; grid coarsening is not an independent proof that
the underlying curve is calibrated. Computational tail resolution and
sampling error are distinct.

A two-sided difference calculation uses

\[
\bar p_{\mathrm{diff}}(\theta_0)=
\min\{1,2\min(\bar p_+(\theta_0),\bar p_-(\theta_0))\}.
\]

For a test at level \(a\), include \(a/2\) in \(\mathcal A\). The
primary practical test uses
\(H_0:\theta_{\mathrm{primary}}\leq\delta_{\mathrm{feedback}}\) at the
proposed one-sided level 0.025, equivalently the strict
lower-95\%-interval rule. Numerical and sampling-calibration
requirements apply to it as to the secondary calculations. No calibrated
test or rejection is currently available because no independent pilot
blocks or formal outcomes exist.

\subsection{Difference, equivalence, and
multiplicity}\label{difference-equivalence-and-multiplicity}

Preserve the three-member secondary difference family:
\(\theta_1=I_{\mathrm{three}}\), \(\theta_2=I_{\mathrm{position}}\), and
\(\theta_3=Q_{11}^{\mathrm{local}}-Q_{11}^{\mathrm{central}}\). Apply
the two-sided calculation at zero to each. For sorted marginal upper
bounds \(\bar p_{(1)}\leq\bar p_{(2)}\leq\bar p_{(3)}\), Holm's
step-down adjusted values are

\[
\bar p^{\mathrm{adj}}_{(j)}=
\min\{1,\max_{i\leq j}(3-i+1)\bar p_{(i)}\}.
\]

Reject sorted hypotheses sequentially while
\(\bar p_{(j)}\leq\alpha_D/(3-j+1)\), stopping at the first failure.
Holm's theorem controls familywise error under arbitrary dependence when
the marginal \(p\) values are valid (Holm 1979, secs. 1--3). Here its
guarantee remains conditional on adequate calibration of the proposed
marginal confidence curves. Applying Holm cannot repair invalid
small-sample inputs. The nominal BCa intervals are marginal intervals,
not Holm-adjusted simultaneous intervals.

For each of these three location contrasts, the draft equivalence bounds
are \(\ell=-0.02\) and \(u=0.02\). TOST tests the union null
\(\theta\leq\ell\) or \(\theta\geq u\), using

\[
\bar p_{\mathrm{eq}}=\max\{\bar p_+(\ell),\bar p_-(u)\}.
\]

Both one-sided nulls must be rejected. At a common marginal level \(a\),
inversion corresponds to \(C_{1-2a}\) lying strictly inside
\((\ell,u)\); touching a bound does not establish equivalence. Thus
each-side 0.05 corresponds to a nominal 90\% interval, whereas a
complete 95\% interval inside the bounds is the stricter each-side 0.025
rule. This test/interval relationship follows the TOST construction
(Schuirmann 1987, 657--62 and 675). Its original normal-theory setting
supplies no finite-sample guarantee for our block-bootstrap application.

An additional three-member equivalence family is proposed, not yet
approved as confirmatory. If approved, apply Holm to its three
\(\bar p_{\mathrm{eq}}\) values separately, with budget \(\alpha_E\).
The recommended joint secondary policy allocates a predeclared secondary
error budget \(\alpha_S\) equally: \(\alpha_D=\alpha_E=\alpha_S/2\). For
an illustrative \(\alpha_S=0.05\), each family receives 0.025. These are
proposed allocation values, not registered thresholds. With valid
marginal tests, the union bound then limits the error probability across
the six secondary claims by \(\alpha_D+\alpha_E\). Running both families
at 0.05 would only establish two separate family controls. This proposal
does not give joint control over the primary claim plus all secondary
claims.

Register the exact families, error budgets, grid levels for every Holm
threshold, and decision rules before formal data. A nominal 90\%
interval inside the equivalence range does not bypass family adjustment.
Non-rejection of a difference is not equivalence; either decision can
remain inconclusive. L11--L10, L11--L01, perturbations, and process
analyses are descriptive unless separately designated and budgeted
before collection. The ASA press release motivates transparent effect
reporting (American Statistical Association 2016); technical support for
these procedures comes from the original articles cited above.

\subsection{Sample-size planning}\label{sample-size-planning}

Formal replication counts depend on pilot estimates of the
sequence-level interaction variance, task dependence, the minimum useful
effect, and budget. The supplied planning example considers a true
interaction of 0.06 against a target of 0.03 at 80\% power. Under an
approximate independent-block normal model, a planning calculation has
the form

\[
r\ \approx\left(\frac{(z_{0.975}+z_{0.80})\sigma_I}{0.06-0.03}\right)^2,
\]

where \(\sigma_I\) must come from independent pilot block contrasts. The
normal approximation concerns independent blocks and the proposed
primary 97.5\% one-sided bound; it does not prove power for the BCa
procedure. Real pilot sequences are needed to assess variance,
dependence, discreteness, degeneracy, and usable precision. Coverage or
null calibration additionally needs justified sampling assumptions and
an independently specified calibration assessment; repeatedly resampling
one empirical pilot distribution cannot establish population coverage by
itself. No such assessment, pilot variance, or numerical replication
count exists here. More Monte Carlo resamples reduce computational error
without adding independent experimental blocks. Replication counts,
\(B\), maximum runs, stopping rules, and numerical diagnostics remain
open. Formal outcomes must not determine them retrospectively.

Propose the following readiness and decision order before confirmatory
collection:

\begin{enumerate}
\def\labelenumi{\arabic{enumi}.}
\tightlist
\item
  Establish lawful fixed public-project inputs, actual native-controller
  tool restrictions, current grants, source access and externally
  directed actions. Resolve the pending S3/native choices. A development
  model-call allowance does not satisfy those prerequisites.
\item
  Implement and independently check the common research policy, eight
  isolated states, both update projections, actual decision views and
  source-linked verdict semantics. Verify evaluation read denial and
  leakage paths with actual authorized records. A method specification
  or role matrix does not complete this step.
\item
  Demonstrate auditable usage, retry/inclusion reconciliation,
  preparation/restoration sources and current budget enforcement in the
  approved units. Preserve unknowns. Before observing pilot outcomes,
  version its models, prompts, inference settings, tools,
  experience/relationship update policy, task distribution and sequence
  length \(T\), initial histories, budgets, permissions, planned block
  count and stopping/exclusion rules. Identify the intended formal
  target population and contrast whose variance the pilot is meant to
  inform. These are a prospective pilot specification and planning
  target, not choices approved by this manuscript.
\item
  Execute independent complete eight-configuration pilot blocks with
  separate condition requests and isolated state. Pilot development
  feedback and evaluation are kept outside the formal outcome set.
  Restoring one historical snapshot can support conditional pilot
  replication, but does not independently sample historical populations.
  Record the task/history population and any provider or cross-block
  dependence.
\item
  Use the actual complete-block contrasts to assess variance
  uncertainty, ties, distinct values, dependence, zero-variance cases,
  BCa diagnostics and usable tail/Monte Carlo precision. A zero or
  unstable pilot variance does not justify a zero or arbitrarily small
  formal sample. Assess the marginal confidence curve's calibration
  under justified assumptions and an independently specified assessment;
  pilot resampling alone supplies no population coverage or power
  guarantee. If applicability or calibration remains unsupported, formal
  inference remains unready.
\item
  Only then confirm formal tasks/models, practical and equivalence
  thresholds, budgets, \(r\), \(B\), \(J\), grid and tail-rank criteria,
  numerical tolerances, randomization, maximum runs, stopping/exclusions
  and sealed-release rules. Compare the formal target against the
  versioned pilot specification component by component. A substantive
  change in models, prompts, inference/tools, update policy, task
  population/\(T\), history or budgets cannot inherit the old
  \(\sigma_I\) without an explicit scientific applicability argument. If
  that argument is unavailable, preplan independent complete pilot
  blocks for the revised target before using its variance to choose
  \(r\). Additional pilot blocks remain development data and costs,
  outside the formal block set. Document the uncertainty and target
  behind the chosen \(r\), freeze the analysis and obtain the required
  decisions before collecting any formal block. The normal planning
  illustration does not replace a justified power/precision assessment
  for the chosen procedure.
\end{enumerate}

Disjoint pilot and formal tasks need a justified population
relationship; different task lists do not automatically estimate the
same conditional variance. If formal inference is conditional on a
particular fixed list/history, variance from another list/history needs
a corresponding applicability argument. A planning revision never
authorizes looking at formal outcomes or moving pilot blocks into the
confirmatory sample. None of the target versions, additional pilot runs
or formal parameters is selected or executed here.

No independent pilot or calibration assessment has been run here, and no
numerical \(r\) is available. Extra tasks within a block, requested
seeds, diagnostic batches and repeated analyses do not increase the
independent block count. Formal scores cannot be inspected to adjust
thresholds, add favorable blocks, choose a grid, or revise the
allocation denominator. If prerequisites or required settings remain
undecided, the study remains at preparation; it does not acquire an
implicit default from this draft.

\section{Adaptation, mechanisms, and prospective iteration
case}\label{adaptation-mechanisms-and-prospective-iteration-case}

\subsection{Member perturbations}\label{member-perturbations}

The protocol proposes member exit, entry, announced capability change,
and unannounced capability change, comparing L11, L00, and C11.
Perturbation times, proportions, identities, and capabilities remain
scenario-registration fields. Matched configurations begin from
corresponding states and schedules. An exit updates availability
immediately; its effect on ongoing work requires real lifecycle
evidence.

Measures include task transfer, incorrect delegation, detection delay,
relationship adjustment, recovery time, quality loss, and cost.
Selection records reveal whether a new member receives opportunities and
whether reused experience supports subsequent work. A stopped Mission
with no Worker provides no evidence about adaptation to member exit.

\subsection{Structural distance and negative
interactions}\label{structural-distance-and-negative-interactions}

Task distance is computed before execution relative to the initial
historical tasks. The draft uses normalized clipped numerical
differences, categorical disagreement, and set differences based on
intersection and union, averaged over observed features. Insufficiently
observed tasks form an unknown group. The exact scale, nearest-history
rule, and near/middle/far boundaries must be fixed using development
data.

Distance groups share the same tasks across conditions. Group-level
L11--L00 gains are reported with coverage and difficulty indicators.
Distance is observed rather than randomized, so a gain-distance pattern
does not independently establish a causal transfer law.

Process records support examination of harmful experience reuse,
concentration of assignments, stale capability versions, delayed
updates, and learning overhead. These are candidate explanations.
Attribution requires separately labeled interventions, such as disabling
a particular experience or changing update visibility in a paired rerun.
Replaying decisions cannot reproduce unobserved model or tool outcomes
in a different environment. Observational logs alone cannot establish
the cause of a negative interaction.

\subsection{Swarm A/B/C application
case}\label{swarm-abc-application-case}

The proposed case assigns hypothesis formation to A, candidate
modification to B, and development verification and independent
evaluation to C. Each candidate records a prediction, applicability
conditions, code/configuration version, development checks, adoption
decision, cost, and lineage. These are functional responsibilities; this
manuscript does not claim that autonomous A/B/C iteration has occurred.

Three processes compare persistent cross-round experience, reset
cross-round experience, and an unchanged system under matched budgets.
Development feedback supports candidate selection. Sealed evaluation is
controlled by the evaluator; its results must not influence later
candidate generation during the evaluated sequence of rounds. Repeated
access to the same final holdout would compromise that separation.

Round-level quality, transfer to new tasks, adopted-improvement cost,
and cumulative performance will be reported when measured. A
hypothesis-acceptance rate describes the research process and cannot
substitute for independent performance. All candidates, failures, and
costs must remain in the case record. This first draft contains no
invented round table.

\section{Current implementation and available
evidence}\label{current-implementation-and-available-evidence}

The inspected product candidate is
\texttt{7ce3a5efe312f4cdac3491293e002f5559f076b0}; its business-code
source is \texttt{982333b871520e606c74d553e448c5049f81abde}. The
research writing branch began from a separate repository default and
does not imply an updated product checkout. The following evidence
levels are kept distinct.

{\def\LTcaptype{none} % do not increment counter
\begin{longtable}[]{@{}
  >{\raggedright\arraybackslash}p{(\linewidth - 4\tabcolsep) * \real{0.3333}}
  >{\raggedright\arraybackslash}p{(\linewidth - 4\tabcolsep) * \real{0.3333}}
  >{\raggedright\arraybackslash}p{(\linewidth - 4\tabcolsep) * \real{0.3333}}@{}}
\toprule\noalign{}
\begin{minipage}[b]{\linewidth}\raggedright
Level
\end{minipage} & \begin{minipage}[b]{\linewidth}\raggedright
Current scope
\end{minipage} & \begin{minipage}[b]{\linewidth}\raggedright
Implication for this manuscript
\end{minipage} \\
\midrule\noalign{}
\endhead
\bottomrule\noalign{}
\endlastfoot
Proposed method & Conditional experience, relationship updates, common
probabilistic policy, eight configurations & Specified here; complete
execution remains unverified \\
Inspected implementation & Mission/WorkItem and lease logic; fixed
candidate ranking; human-governed experience; accepted-source linkage &
Reusable product components with narrower semantics \\
Independent engineering report & Selected manually created DAG,
permission, reconnect, lifecycle-without-Worker, and loading checks &
Component evidence reported by a separate agent session \\
Formal research results & Eight-configuration study, member adaptation,
autonomous iteration & Not collected \\
\end{longtable}
}

In \texttt{internal/swarm/domain/policy.go}, \texttt{swarm.fixed.v1}
filters by approval, budget, capabilities, availability of dependencies,
and write-path conflicts, then orders by matching inputs and identifier.
It does not establish the proposed experience-and-relationship softmax
policy. The inspected \texttt{internal/attention/domain/experience.go}
and application service distinguish quarantine and approved entries with
human governance. Accepted-artifact linkage exists in
\texttt{internal/swarm/app/experience.go}. These findings support an
implementation plan; they do not demonstrate realized experience-driven
collective adaptation.

The independent report is pinned separately (\emph{Astrocyte Dynamic
{DAG} Independent Engineering Review} 2026). Its accepted scenarios use
manually created work items and do not validate model-generated task
planning or Worker exit. This research worker read that report and
inspected the stated source; it did not rerun the product tests or
independently query their event database.

Restricted native execution remains blocked by the requirement to
disable native file and command tools. The pending S3 decisions concern
call-budget accounting and the evidence required for tools and external
destinations; S3 here names the project's Swarm execution stage. They
must be resolved for the affected research execution path. Unknown
delivery cannot be retried as though unsent. OpenViking's positive
remote path, actual work completion/adoption, and subsequent experience
reuse remain unverified.

Full CI has not passed. The reported automated runs were ultimately
cancelled, and their Ubuntu full-repository Go test jobs had already
failed before cancellation. Some earlier fixture and migration
inconsistencies were identified; failure attribution was incomplete.
Limited component acceptance cannot establish whole-product correctness.

The separately governed V4 study is outside this protocol. Its original
external evidence root is unavailable according to the current project
handoff. Historical aggregates without their inspectable raw evidence
are excluded from the present research results. No V4 configuration,
lock, database, score, or automation was changed for this draft.

\section{Results status and interpretation
limits}\label{results-status-and-interpretation-limits}

No main-study, perturbation, or iteration-case results have been
collected for this manuscript. There is no estimated \(Q\), update
interaction, interval, significance value, equivalence conclusion, cost
improvement, or measured adaptation curve to report. Implementation
inspection and literature reading establish readiness information only.

When data become available, reporting must include all eight condition
outcomes, paired contrasts, uncertainty, failure denominators, cost
components, and the actual sampling population. A beneficial total
update gain with negligible interaction would support useful additive
effects. A negative interaction would motivate controlled mechanism
investigation. A near-zero location contrast with sufficiently precise
intervals could describe similar behavior under the tested deployment; a
wide interval would remain inconclusive. None of these outcomes is
selected in advance.

\section{Limitations, reproducibility, and research
governance}\label{limitations-reproducibility-and-research-governance}

The proposed semantics have not been implemented as a complete
executable research system. The concrete condition catalogue,
unique-verdict rule and clock are proposals for review, without formal
parameter or interface approval. Conditions may omit relevant task
differences, and partner selection changes which relationships are
observed. Calibration among selected partners cannot automatically
generalize to unselected pairs. Fixed software tasks and one initial
history limit external validity. Provider-side randomness and
model-version changes can weaken paired comparisons even when requested
seeds match.

The evidence-quality defaults do not guarantee broad applicability.
Feedback tests may be exploitable or overlap with final evaluation;
human ratings may vary; model agreement can share errors. Authorization
and privacy remain constraints on available evidence, with no permission
inferred from the existence of an experience store. Initial snapshots
and the formal task corpus require lawful, explicit provenance.

Before registration, the reproducibility package must identify the exact
code and policy versions, all model/service versions, prompts and
decoding settings, task sequences and suite access, initial-state
sources, per-condition isolation, hardware and communication placement,
resource limits, costs, randomization, exclusions, replication counts,
and analysis programs. Existing project services and records should be
reused. This requirement does not authorize a new scheduler or
proof-of-completion system.

Formal execution depends on legitimate native-controller configuration,
real persistent state and HTTP services, and the actual session files
produced by the controller. Each claimed event must be locatable in the
real event database. The writing work reported here sends no new
product-native model input and runs no model experiment. Development
writing tools are separate from permissions for the proposed
experimental agents.

The supplied protocol remains a draft. Authors, affiliations, ORCID
identifiers, final thresholds, budgets, tasks, replication counts, and a
registration identifier are unfilled. Public submission and formal
registration require subsequent decisions. The manuscript is intended
for scientific and implementation review in its present evidence state.

\subsection{Use of AI tools}\label{use-of-ai-tools}

Codex CLI 0.162.0 using the Orca-configured \texttt{gpt-6.1-sol} model
assisted source retrieval, synthesis, manuscript editing, and
build/verification script preparation. No experimental model outcomes
were generated for this manuscript. Human authorship, scientific
verification, approval, and responsibility for a public submission
remain to be established. This disclosure records actual tool use; it
does not attest that currently unidentified authors have completed those
checks.

\section{Conclusion}\label{conclusion}

Astrocyte specifies a study of coupled external experience and
collaboration-relationship updates under fixed model parameters. The
eight-configuration design distinguishes update interaction,
decision-location modification of interaction, and decision-location
modification of total gain. Its evidential requirements include paired
complete sequences, isolated adaptive states, protected evaluation, and
comprehensive accounting. Current product evidence supports selected
components; the proposed research effects remain to be measured.

\section{References}\label{references}

\protect\phantomsection\label{refs}
\begin{CSLReferences}{1}{1}
\bibitem[\citeproctext]{ref-asa2016release}
American Statistical Association. 2016. \emph{American Statistical
Association Releases Statement on Statistical Significance and
p-Values}.
\url{https://www.amstat.org/asa/files/pdfs/p-valuestatement.pdf}.

\bibitem[\citeproctext]{ref-astrocyteEngineering2026}
\emph{Astrocyte Dynamic {DAG} Independent Engineering Review}. 2026.
\url{https://github.com/songconmaisaix31-design/Astrocyte/blob/7f781deea143c74bed5303f90d8130c82bb54bca/docs/reviews/dynamic-dag-20261019.md}.

\bibitem[\citeproctext]{ref-diciccio1996}
DiCiccio, Thomas J., and Bradley Efron. 1996. {``Bootstrap Confidence
Intervals.''} \emph{Statistical Science} 11 (3): 189--228.
\url{https://doi.org/10.1214/ss/1032280214}.

\bibitem[\citeproctext]{ref-evomapSwarm}
EvoMap. n.d. \emph{Swarm Intelligence}. Accessed October 19, 2026.
\url{https://evomap.ai/wiki/10-swarm}.

\bibitem[\citeproctext]{ref-hall1991}
Hall, Peter, and Susan R. Wilson. 1991. {``Two Guidelines for Bootstrap
Hypothesis Testing.''} \emph{Biometrics} 47 (2): 757--62.
\url{https://ruc-econometrics.weebly.com/uploads/2/3/3/8/23387606/two_guidelines_for_bootstrap_hypothesis_testing.pdf}.

\bibitem[\citeproctext]{ref-hao2026decentmem}
Hao, Guangya, Yunbo Long, and Zhuokai Zhao. 2026. \emph{Self-Evolving
Multi-Agent Systems via Decentralized Memory}.
\url{https://arxiv.org/html/2605.22721v1}.

\bibitem[\citeproctext]{ref-holm1979}
Holm, Sture. 1979. {``A Simple Sequentially Rejective Multiple Test
Procedure.''} \emph{Scandinavian Journal of Statistics} 6 (2): 65--70.
\url{https://www.ime.usp.br/~abe/lista/pdf4R8xPVzCnX.pdf}.

\bibitem[\citeproctext]{ref-hu2026sharemem}
Hu, Jinming, Haodong Zhao, Qi Jia, et al. 2026. \emph{Learning from
Others, Acting for You: Cross-User Memory Sharing for {LLM} Agents}.
\url{https://arxiv.org/html/2609.32511v1}.

\bibitem[\citeproctext]{ref-qu2026maagl}
Qu, Liang, Jianxin Li, Hongzhi Yin, and Hua Wang. 2026.
\emph{Multi-Agent Agentic Graph Learning via Structural Signatures}.
\url{https://arxiv.org/html/2609.09565v1}.

\bibitem[\citeproctext]{ref-schuirmann1987}
Schuirmann, Donald J. 1987. {``A Comparison of the Two One-Sided Tests
Procedure and the Power Approach for Assessing the Equivalence of
Average Bioavailability.''} \emph{Journal of Pharmacokinetics and
Biopharmaceutics} 15 (6): 657--80.
\url{https://doi.org/10.1007/BF01068419}.

\bibitem[\citeproctext]{ref-wang2026genes}
Wang, Junjie, Yiming Ren, and Haoyang Zhang. 2026. \emph{From Procedural
Skills to Strategy Genes: Towards Experience-Driven Test-Time
Evolution}. \url{https://arxiv.org/html/2604.15097v4}.

\bibitem[\citeproctext]{ref-xu2026mpt}
Xu, Guowei, Zhenting Qi, Amani Kiruga, Himabindu Lakkaraju, and Yilun
Du. 2026. \emph{Learning Self-Orchestration in Decentralized Agent
Swarms}. \url{https://arxiv.org/html/2610.22322v1}.

\bibitem[\citeproctext]{ref-zhang2025dgm}
Zhang, Jenny, Shengran Hu, Cong Lu, Robert Lange, and Jeff Clune. 2025.
\emph{Darwin {Gödel} Machine: Open-Ended Evolution of Self-Improving
Agents}. \url{https://arxiv.org/html/2505.22954v1}.

\bibitem[\citeproctext]{ref-zhu2026edv}
Zhu, Shiding, Yudi Qi, Yajie Wang, et al. 2026. \emph{Escaping the
Self-Confirmation Trap: An Execute-Distill-Verify Paradigm for Agentic
Experience Learning}. \url{https://arxiv.org/html/2606.24428v1}.

\end{CSLReferences}

\end{document}
